% Preserved, verbatim, from the original complete Section 6.1 in the 2026-10-09 baseline.
% Source: main.tex blob dc2e6dfd7d34bea5f705cea8814bfe7e148ca731
% Scientific purpose: future migration to Section 6.3 Candidate-Regime Dependence.
% This is not input by main.tex. Do not treat this file as a standalone compiled paper.
% Source evidence: KBS_KUAILIVE_FACTORIAL_SEED_RESULTS_2026-10-08.md
% IMPORTANT: retain post-hoc status, checkpoint identity and P0/P1 distinctions on migration.

\subsection{Candidate-Regime Sensitivity of Base-Relative Memory Utility}
\label{sec:results-regimes}

The KuaiLive sampled-active and full-active evaluations yield different rankings of the fixed relationship-memory specialist relative to the Dual-ID base. To separate candidate-protocol variation from differences between historically trained base configurations, we first evaluate identical user--time--target events with a single fixed Room-SASRec/Streamer-SASRec checkpoint pair and one development-selected fusion coefficient ($\alpha_{\mathrm{room}}=0.125$). For the 10,222 paired test events, Memory--Base NDCG@10 is $-0.04120$ under sampled-active ranking (Base $0.60709$, Memory $0.56589$) and $+0.04631$ under full-active ranking (Base $0.38571$, Memory $0.43202$). The mean full-minus-sampled contrast is $+0.08751$ (95\% user-paired bootstrap interval $[+0.08210,+0.09326]$ in the factorial replay). Both absolute scores decrease when more candidates are ranked, but the base decreases more; the observed sign reversal is therefore a change in \emph{relative}, not absolute, specialist utility.

Table~\ref{tab:results-factorial} crosses the ranked candidate set with the reference population used for the two branch-wise score standardizations, while holding raw model scores, fusion weight, user history, target and Memory rule fixed. The original sampled/sampled and full/full cells reproduce the preceding paired contrast. The mixed cells preserve ranking candidates while changing normalization moments, or preserve normalization moments while changing candidates. Their two-order allocation (Eq.~\eqref{eq:factorial-allocations}) assigns $+0.08827$ of the mean shift to the candidate-membership factor (95\% CI $[+0.08287,+0.09397]$) and $-0.00076$ to the normalization-reference factor (95\% CI $[-0.00117,-0.00036]$). The allocations sum exactly to $+0.08751$. Thus, in this specified evaluation, the change in the \emph{ranked candidate universe} dominates the standardized-score reference change. These are descriptive, path-averaged algebraic allocations rather than independently identified causal effects.

\begin{table}[t]
\centering
\caption{Memory--Base NDCG@10 under a $2\times2$ candidate-membership and score-normalization-reference comparison (10,222 paired KuaiLive test events; checkpoint seed 20260918). Rows specify ranked items; columns specify the z-score reference population.}
\label{tab:results-factorial}
\footnotesize
\setlength{\tabcolsep}{6pt}
\begin{tabular}{@{}lrr@{}}
\toprule
Ranked candidates & Sampled reference & Full-active reference \\
\midrule
Sampled-active & $-0.04120$ & $-0.04142$ \\
Full-active & $+0.04761$ & $+0.04631$ \\
\bottomrule
\end{tabular}
\end{table}

The direction of the candidate-regime contrast is robust across the three independently trained checkpoint pairs (Table~\ref{tab:results-seeds}). Sampled-active Memory--Base is negative and full-active Memory--Base is positive for all three seeds. The mean paired shift across seeds is $+0.08508$, with a between-seed sample standard deviation of $0.00212$ (range $+0.08363$ to $+0.08751$). The mean candidate-membership allocation is $+0.08563$, compared with $-0.00056$ for normalization-reference changes. The three training realizations support directional stability within this model configuration but are insufficient to characterize a wider seed distribution; each seed's user-bootstrap interval and the across-seed standard deviation describe different sources of variability.

\begin{table}[t]
\centering
\caption{Candidate-regime shifts across independently trained Dual-ID checkpoint pairs. All use identical input data, Memory definitions, sampled/full evaluation targets and room fusion coefficient $0.125$. Differences are NDCG@10 units.}
\label{tab:results-seeds}
\footnotesize
\setlength{\tabcolsep}{6pt}
\begin{tabular}{@{}lrrr@{}}
\toprule
Training seed & Sampled Memory--Base & Full Memory--Base & Full$-$Sampled \\
\midrule
20260918 & $-0.04120$ & $+0.04631$ & $+0.08751$ \\
20260919 & $-0.05157$ & $+0.03206$ & $+0.08363$ \\
20260920 & $-0.04904$ & $+0.03505$ & $+0.08408$ \\
\bottomrule
\end{tabular}
\end{table}

Evidence-state prevalence is unchanged by the paired candidate swap: represented, recoverable-but-unrepresented and unavailable states account for 4,589, 206 and 5,427 users, respectively. Across the three seeds, their average full-minus-sampled Memory--Base shifts are $+0.08618$, $-0.10997$ and $+0.09155$. Their prevalence-weighted contributions to the average overall shift are $+0.03869$, $-0.00222$ and $+0.04860$, summing to $+0.08508$. Consequently, the aggregate reversal does not arise from increased prevalence of recoverable histories, nor from an increased conditional benefit in that small group. It is associated chiefly with increased represented-state relative utility and reduced, though generally still negative, unavailable-state penalties. The recoverable group comprises only 206 events, warranting particular caution in its interpretation.

These comparisons isolate trained-checkpoint and fusion-parameter identities, but neither candidate-set contrasts nor their factor allocations identify a causal effect on user behavior or intrinsic history value. The factorial and seed analyses reuse an already examined test population and are therefore reported as post-hoc supporting evidence; they are not independent held-out validation of the selective policy. Details of model identity, per-user paired outcomes and robustness calculations are retained in the accompanying reproducibility records.

